\section{Conclusion}

Tool-effect binding is the representation obligation that precedes runtime
authorization.  Source execution establishes what a tool commits; a policy
witness establishes which effects require different decisions; and a
representation collision shows whether the monitor preserves that distinction.
When it does not, downstream enforcement cannot recover the lost information.
Source-executed counterfactuals expose these collisions and guide descriptor
refinement before registration.

Across finite AgentDojo and pre-registered ToolSandbox domains, typed effects
remove collisions left by common call views; on ToolSandbox, concrete atoms
exactly realize the specified finite policy relation.  In AgentDojo, a
provenance-origin monitor uses the registered descriptor for deterministic
pre-commit checks and complete call reconciliation.  Its security gains remain
subject to provenance quality and a visible utility tradeoff.  The contribution
is an executable and falsifiable effect representation, not an authority source:
the planner proposes calls, while an independent policy must decide which
represented effects may commit.
\label{technical-body-end}
